\chapter{Bietti's Crystalline Dystrophy}
\index{Bietti's Crystalline Dystrophy}
\label{sec:Bietti's Crystalline Dystrophy}

\section{Overview}
\begin{itemize}[noitemsep]
  \item Bietti's crystalline dystrophy (BCD) is a rare inherited eye disease causing progressive degeneration of the retina and accumulation of crystal-like deposits
  \item It is caused by mutations in the CYP4V2 gene, which plays a role in fat (lipid) metabolism
  \item It is more common in people of East Asian descent and typically begins in young adulthood
\end{itemize}
\section{Causes}
This condition happens because of mutations in the CYP4V2 gene, inherited in an autosomal recessive pattern. The gene normally helps the body process certain fatty acids; when it is defective, lipids accumulate abnormally in the retina and cornea, forming crystalline deposits and causing progressive death of retinal cells, particularly in the pigment epithelium and choroid.
\section{Risk Factors}
\begin{itemize}[noitemsep]
  \item Having two copies of a CYP4V2 mutation (one from each parent)
  \item East Asian (especially Chinese and Japanese) ancestry
  \item Family history of the condition
  \item Consanguineous parents increase the risk of recessive disease
\end{itemize}
\section{Signs and Symptoms}
\subsection{Common symptoms}
\begin{itemize}[noitemsep]
  \item Night blindness (difficulty seeing in dim light)
  \item Progressive loss of peripheral (side) vision
  \item Blurred central vision in later stages
  \item Glare and light sensitivity
  \item Crystals visible in the retina on eye examination, sometimes with similar deposits in the cornea
\end{itemize}
\subsection{Emergency warning signs}
\begin{itemize}[noitemsep]
  \item Rapidly worsening vision requires prompt evaluation to document progression and rule out treatable causes
\end{itemize}
\section{Progression and Stages}
Symptoms usually begin between ages 20 and 40 and progress slowly over decades. Night blindness and peripheral field loss appear first, followed by gradual constriction of the visual fields. Central vision is often preserved until later stages, but eventually significant visual impairment develops. The rate of progression varies between individuals.
\section{Complications}
\begin{itemize}[noitemsep]
  \item Severe visual field constriction (tunnel vision)
  \item Central vision loss in advanced disease
  \item Difficulty driving, reading, and working in dim light
  \item Impaired mobility and independence
  \item Psychological impact and reduced quality of life
\end{itemize}
\section{Diagnosis}
\begin{itemize}[noitemsep]
  \item Dilated fundus examination showing the characteristic retinal crystalline deposits and pigmentary changes
  \item Optical coherence tomography (OCT) showing retinal atrophy
  \item Electroretinography (ERG) to assess retinal function
  \item Visual field testing (perimetry)
  \item Corneal examination for crystals (at the limbus)
  \item Genetic testing to confirm CYP4V2 mutations
\end{itemize}
\section{Prevention}
\begin{itemize}[noitemsep]
  \item Genetic counseling for affected families
  \item Carrier testing and prenatal testing where the family mutation is known
  \item Regular monitoring to detect and treat complications
  \item Use of UV-protective eyewear
  \item Avoidance of smoking
\end{itemize}
\section{Treatment Options}
\begin{itemize}[noitemsep]
  \item There is currently no cure for Bietti's crystalline dystrophy
  \item Gene replacement therapy is under investigation in clinical trials
  \item Low-vision aids and rehabilitation
  \item Vitamin and dietary supplements are sometimes suggested but their benefit is unproven
  \item Treatment of any coexisting conditions such as cataracts
  \item Participation in research registries and trials
\end{itemize}
\section{Living with It}
\begin{itemize}[noitemsep]
  \item Work with low-vision services for magnification and mobility training
  \item Use assistive technology such as screen readers and talking devices
  \item Plan activities around remaining vision, especially for night travel
  \item Join support networks for inherited retinal disease
  \item Stay informed about advances in gene therapy research
\end{itemize}
\section{Outlook}
Bietti's crystalline dystrophy is progressive and currently incurable, and it gradually reduces vision over decades. The pace varies, and many people retain useful vision well into middle age. Research into gene therapy offers cautious hope for slowing or halting the disease in the future.
